% 第 1 章 移动语义的力量（由 tools/convert.py 从原 content/ 自动生成，勿手改）

本章通过一个简短的代码示例，演示了移动语义的基本原理。\par


\section{设计移动语义的动机}

为了了解移动语义的基本原则，使用小段代码为例，首先没有移动语义(例如，用一个旧的C++编译器编译，仅支持C++03)，之后再用移动语义(用现代C++编译器编译，支持C++11或更高的版本)。\par

\mySmallsection{1.1.1 C++03的示例(使用移动语义之前)}

假设我们有以下代码:\par

\begin{codeboxt}[cpp]{basics/motiv03.cpp}
#include <string>
#include <vector>

std::vector<std::string> createAndInsert()
{
  std::vector<std::string> coll; // create vector of strings
  coll.reserve(3); // reserve memory for 3 elements
  std::string s = "data"; // create string object

  coll.push_back(s); // insert string object
  coll.push_back(s+s); // insert temporary string
  coll.push_back(s); // insert string

  return coll; // return vector of strings
}

int main()
{
  std::vector<std::string> v; // create empty vector of strings
  ...
  v = createAndInsert(); // assign returned vector of strings
  ...
}
\end{codeboxt}

使用不支持移动语义的C++编译器编译这段代码后，让我们看一下程序执行的各个步骤(检查堆栈和堆)。\par

\begin{itemize}
  \item 首先，在main()中，我们创建了空的vector \textit{v}:\par
  \textit{\textbf{std::\allowbreak{}vector<\allowbreak{}std::\allowbreak{}string> v;}} \par
  存放在堆栈上，元素数为0，未为元素分配内存。
  \item 然后，调用\par
  \textit{\textbf{v =\allowbreak{} createAndInsert(\allowbreak{})\allowbreak{};}}\par
  在堆栈上创建另一个vector \textit{coll}，并在堆上为三个元素预留内存:\par
  \textit{\textbf{std::\allowbreak{}vector<\allowbreak{}std::\allowbreak{}string> coll;}}\par
  \textit{\textbf{coll.reserve(\allowbreak{}3)\allowbreak{};}}\par
  因为分配的内存没有初始化，所以元素的数量仍然是0。
  \item 然后，用“data”初始化字符串:\par
  \textit{\textbf{std::\allowbreak{}string s =\allowbreak{} "data";}}\par
  字符串类似于具有\textit{char}元素的vector。其实是在堆栈上创建一个对象，其中一个成员表示字符的数量(值为4)，另一个指针表示字符的内存。\par
  之后，程序的状态如下:堆栈上有三个对象:\textit{v}、\textit{coll}和\textit{s}。其中\textit{coll}和\textit{s}已经在堆上分配了内存:\par
\myGraphic{0.8}{images/chapter1/1}{}
  \item 下一步是将字符串放到\textit{coll}中:\par
  \textbf{\textit{coll.push\_\allowbreak{}back(\allowbreak{}s)\allowbreak{};}}\par
  C++标准库中的所有容器都具有语义，能创建副本。所以，vector有了第一个元素，它是\textit{s}的副本:\par
\myGraphic{0.8}{images/chapter1/2}{}
  目前为止，我们还没有对这个程序进行优化。现在，我们有两个vector，\textit{v}和\textit{coll}，以及两个字符串，\textit{s}和它的副本，这也是\textit{coll}的第一个元素。它们都是独立的对象，有自己的内存。\par
  \item 现在让我们看看下一句，创建了一个新的临时字符串，并再次将其插入vector中:\par
  \textbf{\textit{coll.push\_\allowbreak{}back(\allowbreak{}s+s)\allowbreak{};}}\par
  该语句分三个步骤执行:\par
  \begin{enumerate}
    \item 创建临时字符串s+s:\par
\myGraphic{0.8}{images/chapter1/3}{}
    \item 将这个临时字符串插入到\textit{coll}中。容器创建副本，创建临时字符串的深层副本，包括分配内存:\par
\myGraphic{0.8}{images/chapter1/4}{}
    \item 末尾，临时字符串\textit{s+s}销毁:\par
\myGraphic{0.8}{images/chapter1/5}{}
    这是我们第一次生成代码执行不大合适的地方：创建临时字符串的副本后销毁临时字符串。其实，可以不必分配和释放内存，直接使用原始的临时字符串对象。
  \end{enumerate}
  \item 下一行中，我们再将\textit{s}插入\textit{coll}中:\par
  \textbf{\textit{coll.push\_\allowbreak{}back(\allowbreak{}s)\allowbreak{};}}\par
  \textit{coll}继续拷贝\textit{s}:\par
\myGraphic{0.8}{images/chapter1/6}{}
  这也是需要优化：一些优化技巧可以使用\textit{s}作为vector的新元素。
  \item 在\textit{createAndInsert(\allowbreak{})\allowbreak{}}的末尾，使用\textit{return}:\par
\begin{codebox}[cpp]
    return coll;
  }
\end{codebox}
  程序的行为变得有点复杂。按值返回(返回类型不是引用)的应该是\textit{coll}的副本。创建副本意味着必须创建整个vector，及其所有元素的副本。因此，必须为vector中的元素数组分配堆内存，并为每个字符串分配用于保存其值的值分配堆内存。这样，我们需要分配4次内存。\par
  由于\textit{coll}被销毁，编译器可以执行命名返回值优化(NRVO)。这意味着编译器可以生成代码，使\textit{coll}只是用作返回值。\par
  可以使用这种优化，不过这可能会改变程序的行为。如果在vector或string对象的复制构造函数中有print语句，则会不看print语句的再次输出，这意味着这种优化改变了程序的功能。然而，这是可以的，因为在C++标准中允许这种优化，即使它有副作用。任何人都不应该知道这里有一个副本，是否执行命名返回值优化仅取决于编译器。\par
  假设我们进行了命名返回值优化，在\textit{return}语句的末尾，\textit{coll}现在成为返回值，并调用\textit{s}的析构函数，释放声明\textit{s}时分配的内存:\par
\myGraphic{0.8}{images/chapter1/7}{}
  \item 最后，将返回值赋给\textit{v}:\par
  \textbf{\textit{v =\allowbreak{} createAndInsert(\allowbreak{})\allowbreak{};}}\par
  这里可以改进的行为:通常的赋值操作符的目标是赋予\textit{v}与赋值源的值相同的值。一般来说，任何值都不应该修改，并且应该独立于赋值的对象。因此，赋值操作符将创建了一个返回值的副本:\par
\myGraphic{0.8}{images/chapter1/8}{}
  之后，就不再需要临时返回值了，可以销毁:\par
\myGraphic{0.8}{images/chapter1/9}{}
  我们创建了一个临时对象的副本，然后立即销毁该副本的源，再次分配和释放内存。这次适用于四个分配，一个分配给vector，其余分配给每个string元素。
\end{itemize}

对于\textit{main()}中赋值后的状态，分配了10次内存，释放了6次。不必要的内存分配由以下原因引起:\par

\begin{itemize}
  \item 将临时对象插入容器
  \item 将对象插入到不再需要该值的容器中
  \item 将一个临时vector的值赋于其他对象
\end{itemize}

我们可以或多或少地避免这些性能损失。我们可以这样做:\par

\begin{itemize}
  \item 将vector作为out参数传递:\par
  \textit{\textbf{createAndInsert(v); {\color{mygreen}{// let the function fill vector v}}}}
  \item 使用swap():\par
  \textit{\textbf{createAndInsert(\allowbreak{})\allowbreak{}.swap(\allowbreak{}v)\allowbreak{};}}
\end{itemize}

但是，得到的代码看起来更难看懂(除非您在复杂的代码中看到一些美)，而且在插入临时对象时没有真正的解决问题。\par

自C++11起，我们有了另一个选择:编译并运行支持移动语义的程序。\par

\mySmallsection{1.1.2 使用C++11的例子(使用移动语义)}

现在让我们用一个支持移动语义的现代C++编译器(C++11或更高版本)重新编译这个程序:\par

\begin{codeboxt}[cpp]{basics/motiv11.cpp}
#include <string>
#include <vector>

std::vector<std::string> createAndInsert()
{
  std::vector<std::string> coll; // create vector of strings
  coll.reserve(3); // reserve memory for 3 elements
  std::string s = "data"; // create string object

  coll.push_back(s); // insert string object
  coll.push_back(s+s); // insert temporary string
  coll.push_back(std::move(s)); // insert string (we no longer need the value of s)

  return coll; // return vector of strings
}

int main()
{
  std::vector<std::string> v; // create empty vector of strings
  ...
  v = createAndInsert(); // assign returned vector of strings
  ...
}
\end{codeboxt}

这里有一个小修改:将最后一个元素插入\textit{coll}时，添加了\textit{std::move()}的调用，我们将在讨论这个语句时聊聊这个更改。其他都和之前一样。\par

让我们再次通过检查堆栈和堆来了解程序的各个步骤。\par

\begin{itemize}
  \item 首先，在\textit{main()}中创建空vector \textit{v}，包含0个元素：\par
  \textit{\textbf{std::\allowbreak{}vector<\allowbreak{}std::\allowbreak{}string> v;}}
  \item 之后，调用：\par
  \textit{\textbf{v =\allowbreak{} createAndInsert(\allowbreak{})\allowbreak{};}}\par
  在堆栈上创建另一个空vector \textit{coll}，并保留三个元素未初始化的内存：\par
  \textit{\textbf{std::\allowbreak{}vector<\allowbreak{}std::\allowbreak{}string> coll;}}\par
  \textit{\textbf{coll.reserve(\allowbreak{}3)\allowbreak{};}}\par
  \item 然后，我们创建用“data”初始化的字符串\textit{s}，并再次将其插入\textit{coll}中：\par
  \textit{\textbf{std::\allowbreak{}string s =\allowbreak{} "data";}}
  \textit{\textbf{coll.push\_\allowbreak{}back(\allowbreak{}s)\allowbreak{};}}
  目前为止，还没有什么需要优化的，程序状态与C++03相同：\par
\myGraphic{0.8}{images/chapter1/10}{}
  我们有两个vector，\textit{v}和\textit{coll}，还有两个字符串，\textit{s}和它的副本，也就是\textit{coll}的第一个元素。它们都是独立的对象，有自己的内存。\par
  \item 首先，看看创建一个新的临时字符串并将其插入vector中的语句:\par
  \textit{\textbf{coll.push\_\allowbreak{}back(\allowbreak{}s+s)\allowbreak{};}}\par
  同样，该语句分三个步骤执行:\par
  \begin{enumerate}
    \item 创建临时字符串\textit{s+s}:\par
\myGraphic{0.8}{images/chapter1/11}{}
    \item 将这个临时字符串插入到\textit{coll}中。然而，这里发生了一些不同的事情:我们从\textit{s+s}中窃取了内存，并将其移动到\textit{coll}的新元素中。\par
\myGraphic{0.8}{images/chapter1/12}{}
    自C++11起，我们可以获取不再需要的值。因为编译器知道在执行push\_back()调用之后，临时对象\textit{s+s}将被销毁。因此，调用的是push\_back()的另一个实现，其是对不再需要值的字符串进行复制的优化：我们复制了大小和指向内存的指针，而不是创建拷贝。然而，这种浅复制是不够的;还修改了临时对象\textit{s+s}，将其大小设置为0，并将nullptr赋值为新值。其实，s+s被修改了，以便呈现空字符串的状态。重要的是，它不再拥有自己的内存。
    \item 语句的结尾，临时字符串\textit{s+s}被销毁。但由于临时字符串不再是初始内存的所有者，析构函数不会释放该内存。\par
\myGraphic{0.8}{images/chapter1/13}{}
    我们优化了复制过程，将\textit{s+s}的内存所有权转移到它在vector中的副本。\par
    这一切都是通过编译器自动完成的，编译器可以对即将死亡的对象发出信号，可以使用新实现来复制一个从源端窃取值的字符串值。这不是一个技巧，这是一种语义上的移动，通过从底层技术上将值的内存从源字符串移动到它的副本来实现。\par
  \end{enumerate}
  \item 下一条语句是我们为C++11版本修改的语句。将\textit{s}插入到\textit{coll}中，但是通过对要插入的字符串\textit{s}调用\textit{std::move()}，语句发生了改变:\par
  \textit{\textbf{coll.push\_\allowbreak{}back(\allowbreak{}std::\allowbreak{}move(\allowbreak{}s)\allowbreak{})\allowbreak{};}}\par
  如果没有\textit{std::move()}，就会像第一次调用\textit{push\_\allowbreak{}back(\allowbreak{})\allowbreak{}}一样：vector会创建传递的字符串\textit{s}的副本。这次，我们用\textit{std::move()}标记了\textit{s}，这在意味着“这里不再需要这个值”。因此，这里会使用另一个\textit{push\_\allowbreak{}back(\allowbreak{})\allowbreak{}}实现，会在传递临时对象\textit{s+s}时使用。第三个元素通过将这个值的内存所有权，从\textit{s}转移到副本中，从而来窃取这个值:\par
\myGraphic{0.8}{images/chapter1/14}{}
  关于移动语义，请注意以下两点:\par
  \begin{itemize}
    \item[-] \textit{std::\allowbreak{}move(\allowbreak{}s)\allowbreak{}}在这里只表示\textit{s}是可移动的，只表示不再需要这个值，允许实现通过在复制值时进行一些优化(比如偷取内存)从而获益。调用者并不知道值是否进行了移动。
    \item[-] 然而，窃取值的优化必须确保源对象仍然处于有效状态。被移动的对象既不会部分销毁，也不会完全销毁。C++标准库为它的类型进行了规定:标记为\textit{std::move()}的对象执行操作后，该对象处于有效，但未定义的状态。\par
    以上，就是在这条语句执行完成后会发生的事\par
    \textit{\textbf{coll.push\_\allowbreak{}back(\allowbreak{}std::\allowbreak{}move(\allowbreak{}s)\allowbreak{})\allowbreak{};}}\par
    这里保证\textit{s}仍然是有效的字符串。这就像使用不知道传递了哪个值的字符串参数一样。\par
    注意，它也不能保证字符串要么有旧值，要么为空，由运行时库实现者决定。通常，实现者可以对\textit{std::move()}操作的对象做任何事，只要保持对象的有效状态。保证的理由，稍后讨论。
  \end{itemize}

  \item \textit{createAndInsert(\allowbreak{})\allowbreak{}}最后的return语句为:\par
\begin{codebox}[cpp]
    return coll;
  }
\end{codebox}
  是否生成具有指定返回值优化的代码取决于编译器，\textit{coll}只是成为了返回值。但是，如果不使用这种优化，return语句仍然是从即将失效的源创建对象的情况。如果没有使用指定的返回值优化，仍然会使用移动语义，返回值会从\textit{coll}中窃取值。最坏的情况下，必须从源文件复制成员的大小、容量和指向内存的指针(总共12或24个字节)到返回值中，并在源文件中为这些成员赋值。\par
  我们假设进行了返回值优化。在return语句中，\textit{coll}成为返回值，并调用\textit{s}的析构函数，不再需要释放任何内存，因为内存移到了\textit{coll}的第三个元素中:\par
\myGraphic{0.8}{images/chapter1/15}{}
  \item 最后，来给\textit{v}赋值：\par
  \textit{\textbf{v =\allowbreak{} createAndInsert(\allowbreak{})\allowbreak{};}}\par
  同样可以从移动语义中获益:必须从一个即将销毁的临时返回值中复制(这里是赋值)一个值。\par
  现在，移动语义允许从源vector转移值的赋值操作符:\par
\myGraphic{0.8}{images/chapter1/16}{}
  同样，临时对象不会(部分地)销毁。它进入了某个状态，但我们不知道它的值。\par
  但在赋值之后，语句的末尾会销毁(修改过的)临时返回值:\par
\myGraphic{0.8}{images/chapter1/17}{}
\end{itemize}

最后，和之前使用移动语义时一样，但有一些事情发生了改变:节省了6个内存分配和释放。所有不必要的内存分配不再发生:\par

\begin{itemize}
  \item 将临时对象移动插入容器
  \item 当使用\textit{std::move()}表示不再需要该值时，用于将对象移动插入容器
  \item 临时vector及其元素的分配方式
\end{itemize}

第二种情况下，优化是在我们的帮助下完成的。通过添加\textit{std::move()}，说明不再需要\textit{s}的值。所有其他优化都生效了，因为编译器知道对象即将销毁，所以编译器也可以使用移动语义的优化实现。\par

这意味着返回字符串向量，并将其赋值给现有向量不再有性能问题。我们可以像使用整型一样使用字符串向量，从而获得更好的性能。在实践中，用移动语义重新编译代码可以将速度提高10\%到40\%(取决于现有代码的难易程度)。\par


\section{实现移动语义}

用前面的例子来了解移动语义是如何实现的。\par

实现移动语义之前，\textit{std::\allowbreak{}vector<\allowbreak{}>}只有一个\textit{push\_\allowbreak{}back(\allowbreak{})\allowbreak{}}的实现(这里简化了vector的声明):\par

\begin{codebox}[cpp]
template<typename T>
class vector {
  public:
  ...
  // insert a copy of elem:
  void push_back (const T& elem);
  ...
};
\end{codebox}

传递参数给\textit{push\_\allowbreak{}back(\allowbreak{})\allowbreak{}}的方法只有一种:将其绑定到\textit{const}引用上。\textit{push\_\allowbreak{}back(\allowbreak{})\allowbreak{}}的实现方式是在不修改参数的情况下，创建参数的副本。\par

C++11开始，\textit{push\_\allowbreak{}back(\allowbreak{})\allowbreak{}}有了重载:\par

\begin{codebox}[cpp]
template<typename T>
class vector {
  public:
  ...
  // insert a copy of elem:
  void push_back (const T& elem);
  // insert elem when the value of elem is no longer needed:
  void push_back (T&& elem);
  ...
};
\end{codebox}

第二个\textit{push\_\allowbreak{}back(\allowbreak{})\allowbreak{}}使用了为移动语义引入的新语法。使用两个\&声明实参，不使用\textit{const}，这样的参数称为\textit{\textbf{右值引用}}。只有一个\&的“普通引用”现在称为\textit{\textbf{左值引用}}。在这两个调用中，都传递了通过引用插入的值。但是，区别如下:\par

\begin{itemize}
  \item 使用\textit{push\_\allowbreak{}back(\allowbreak{}const T\&)\allowbreak{}}，可以保证不修改传递的值。\par
  当调用者仍然需要传递的值时，将调用此函数。
  \item 通过\textit{push\_\allowbreak{}back(\allowbreak{}T\&\&)\allowbreak{}}，可以修改传递的实参(因此它不是\textit{const})来“窃取”值。语义上的含义是，新元素接收传递的值，但可以使用优化实现，将值移动到vector中。\par
  当调用者不再需要传递的值时，将调用此函数。实现必须确保传递的参数仍然处于有效状态，但值可以更改。因此，在调用此参数之后，调用者仍然可以使用传递的参数，只要调用者不对其值做任何操作。
\end{itemize}

但\textit{vector}不知道如何复制或移动元素。在确保\textit{vector}有足够的内存容纳新元素之后，\textit{vector}将工作委托给元素的类型。\par

在本例中，元素是字符串。如果我们复制或移动传递的字符串会发生什么。\par

\mySmallsection{1.2.1 使用复制构造函数}

用于传统复制语义的\textit{push\_\allowbreak{}back(\allowbreak{}const T\&)\allowbreak{}}调用string类的复制构造函数，该构造函数初始化vector中的新元素。一个非常简单的string类实现的复制构造函数，看起来是这样的:\par

\begin{codebox}[cpp]
class string {
  private:
  int len; // current number of characters
  char* data; // dynamic array of characters
  public:
  // copy constructor: create a full copy of s:
  string (const string& s)
  : len{s.len} { // copy number of characters
    if (len > 0) { // if not empty
      data = new char[len+1]; // - allocate new memory
      memcpy(data, s.data, len+1); // - and copy the characters
    }
  }
  ...
};
\end{codebox}

假设我们对值为"data"的字符串使用这个复制构造函数:\par

\begin{codebox}[cpp]
std::string a = "data";
std::string b = a; // create b as a copy of a
\end{codebox}

初始化字符串a后，如下所示:\par

\myGraphic{0.9}{images/chapter1/18}{}

上面的复制构造函数将复制成员len以获取字符数，为数据指针分配新的内存，并将源a(作为s传递)中的所有字符复制到新字符串:\par

\myGraphic{0.9}{images/chapter1/19}{}

\mySmallsection{1.2.2 使用移动构造函数}

对于移动语义，\textit{push\_\allowbreak{}back(\allowbreak{}T\&\&)\allowbreak{}}调用相应的构造函数，即\textit{\textbf{移动构造函数}}。从现有字符串创建新字符串的构造函数，和移动语义一样，构造函数使用非\textit{const}右值引用(\textit{\&\&})作为形参:\par

\begin{codebox}[cpp]
class string {
  private:
  int len; // current number of characters
  char* data; // dynamic array of characters
  public:
  ...
  // move constructor: initialize the new string from s (stealing the value):
  string (string&& s)
  : len{s.len}, data{s.data} { // copy number of characters and pointer to memory
    s.data = nullptr; // release the memory for the source value
    s.len = 0; // and adjust number of characters accordingly
  }
  ...
};
\end{codebox}

给定上述拷贝构造函数的情况:\par

\myGraphic{0.9}{images/chapter1/20}{}

我们可以这样调用string对象的构造函数:\par

\begin{codebox}[cpp]
std::string c = std::move(b); // init c with the value of b (no longer needing its value here)
\end{codebox}

移动构造函数首先复制成员len和data，这意味着新字符串获得\textit{b}值的所有权(作为\textit{s}传递)。\par

\myGraphic{0.9}{images/chapter1/21}{}

这是不够的，因为\textit{b}的析构函数将释放内存。因此，我们还可以修改源字符串，使其失去内存的所有权，并使其表示空字符串:\par

\myGraphic{0.9}{images/chapter1/22}{}

结果是，\textit{c}现在有了\textit{b}的值，而且\textit{b}是空字符串。注意，唯一的保证是\textit{b}随后处于有效但未定义的状态。根据C++库中移动构造函数的实现方式，可能不为空(但通常是空的，因为这是提高性能的最简单和最好的方法)。\par


\section{复制是一种应急方式}

可以通过移动语义，移动临时对象或用\textit{std::move()}标记的对象，可以通过“窃取”源值来优化值的复制(采用非\textit{const}右值引用)。但是，如果没有针对移动语义的优化版本，则使用复制。\par

例如，像vector这样的容器类缺少\textit{push\_\allowbreak{}back(\allowbreak{})\allowbreak{}}的重载:\par

\begin{codebox}[cpp]
template<typename T>
class MyVector {
  public:
  ...
  void push_back (const T& elem); // insert a copy of elem
  ... // no other push_back() declared
};
\end{codebox}

仍然可以传递一个临时对象或使用\textit{std::move()}:\par

\begin{codebox}[cpp]
MyVector<std::string> coll;
std::string s{"data"};
...
coll.push_back(std::move(s)); // OK, uses copy semantics
\end{codebox}

对于临时对象或使用\textit{std::move()}的对象，首选是将形参声明为右值引用的函数。如果不存在这样的函数，则使用复制语义。这样，就可以确保调用者不需要知道是否存在优化。优化有可能没有，因为:\par

\begin{itemize}
  \item 函数/类是在移动语义支持之前实现，或者没有考虑支持移动语义
  \item 没有什么需要优化的(只有数字成员的类就是一个例子)
\end{itemize}

对于泛型代码，如果不再需要一个对象的值，可以用\textit{std::move()}。即使没有移动语义的支持，相应的代码也会编译通过。\par

出于同样的原因，甚至可以用\textit{std::move()}来标记基本数据类型的对象，如int(或指针)。仍然会使用复制的值(地址)方式:\par

\begin{codebox}[cpp]
std::vector<int> coll;
int x{42};
...
coll.push_back(std::move(x)); // OK, but copies x (std::move() has no effect)
\end{codebox}


\section{\textit{const}对象的移动语义}

最后，不能移动用\textit{const}声明的对象。因为任何优化实现都要求可以修改传递的实参，如果不允许修改，就不能窃取。\par

使用\textit{push\_\allowbreak{}back(\allowbreak{})\allowbreak{}}的重载:\par

\begin{codebox}[cpp]
template<typename T>
class vector {
  public:
  ...
  // insert a copy of elem:
  void push_back (const T& elem);
  // insert elem when the value of elem is no longer needed:
  void push_back (T&& elem);
  ...
};
\end{codebox}

对\textit{const}对象使用的函数是具有\textit{const\&}形参的\textit{push\_\allowbreak{}back(\allowbreak{})\allowbreak{}}重载:\par

\begin{codebox}[cpp]
std::vector<std::string> coll;
const std::string s{"data"};
...
coll.push_back(std::move(s)); // OK, calls push_back(const std::string&)
\end{codebox}

\textit{const}对象的\textit{std::move()}没起作用。\par

原则上，可以通过声明带有\textit{const}右值引用的函数进行重载，但在语义上没有意义。同样，\textit{const}左值引用会作为处理这种情况的备选。\par

\mySmallsection{1.4.1 \textit{const}返回值}

\textit{const}禁用移动语义对声明返回类型也有影响。\textit{const}返回值不能移动。\par

因此，从C++11开始，用\textit{const}返回值就不再是好的方式了(正如过去的一些风格指南所推荐的那样)。例如:\par

\begin{codebox}[cpp]
const std::string getValue();

std::vector<std::string> coll;
...
coll.push_back(getValue()); // copies (because the return value is const)
\end{codebox}

当按值返回时，不要将整个返回值声明为\textit{const}。仅在声明部分返回类型时使用\textit{const}(例如：返回的引用或指针所指向的对象):\par

\begin{codebox}[cpp]
const std::string getValue(); // BAD: disables move semantics for return values
const std::string& getRef(); // OK
const std::string* getPtr(); // OK
\end{codebox}


\section{总结}

\begin{itemize}
  \item 移动语义允许对对象的复制进行优化，它可以隐式使用(用于未命名的临时对象或局部返回值)，也可以显式使用(通过\textit{std::move()})。
  \item \textit{std::move()}表示不再需要这个值，它将对象标记为可移动的。标记为\textit{std::move()}的对象不会(部分地)销毁(析构函数仍然会调用)。
  \item 通过使用非\textit{const}右值引用(例如\textit{std::\allowbreak{}string\&\&})声明函数，可以定义一个接口，调用者在接口中从语义上声明不再需要传递的值。函数可以通过“窃取”这个信息来进行优化，或者对传递的参数做任何修改。通常，实现者还必须确保传递的参数在调用后处于有效状态。
  \item 移动的C++标准库的对象仍然是有效的对象，但其值为未定义。
  \item 拷贝语义用作移动语义的备选(如果拷贝语义支持的话)。如果没有采用右值引用的实现，则使用任何采用普通\textit{const}左值引用的实现(如：\textit{const std::\allowbreak{}string\&})。即使对象被显式地标记为\textit{std::move()}，也会使用备选方式。
  \item 对\textit{const}对象调用\textit{std::move()}通常没有效果。
  \item 如果按值(而不是按引用)返回，不要将返回值声明为\textit{const}。
\end{itemize}
